% abstract.tex — 摘要

\begin{center}
    \Large\textbf{摘\quad 要}
\end{center}

本文针对2025年全国大学生数学建模竞赛A题，研究了利用无人机投放烟幕干扰弹掩护真目标的多阶段优化决策问题。针对来袭导弹对保护目标的威胁，通过设计无人机的飞行方向、飞行速度、烟幕干扰弹投放点和起爆点等策略参数，实现烟幕遮蔽时间最大化。

\textbf{建模方面}，本文建立了完整的运动学模型，包括导弹直线飞行模型、无人机等高度匀速直线飞行模型、烟幕干扰弹自由落体模型和烟幕云团匀速下沉模型。在此基础上，构建了基于视线(line-of-sight, LOS)遮挡判断的烟幕遮蔽效果评估模型：将真目标圆柱表面离散化为采样点集，通过计算导弹到目标采样点的线段与烟幕球体的空间几何关系，判定遮蔽状态。对多枚干扰弹的遮蔽效果，采用区间合并方法计算总有效遮蔽时长。

\textbf{求解方面}，针对不同复杂度的子问题采用分层递进的求解策略：问题1为确定性仿真计算；问题2采用网格搜索结合序列二次规划(SQP)的混合优化方法；问题3--5采用遗传算法(GA)进行全局优化，并以SQP进行局部精化。针对问题5的高维特征(40个决策变量)，设计了基于UAV-导弹距离的分解预分配策略以降低优化难度。

\textbf{结果方面}，问题1得到确定条件下的有效遮蔽时长为$T_1$秒；问题2通过优化4个决策变量，得到最优遮蔽时长$T_2$秒；问题3的3枚干扰弹策略实现了总遮蔽时长$T_3$秒；问题4的三机协同遮蔽时长为$T_4$秒；问题5在全场景下实现了对各导弹较为均衡的遮蔽效果。灵敏度分析表明，烟幕半径和无人机速度范围对遮蔽效果影响显著。

\textbf{模型优势}在于：(1)几何模型严格，视线遮挡判断基于解析几何，精度可控；(2)分层优化策略兼顾全局搜索能力和局部精度；(3)模块化设计便于扩展到更多无人机和导弹的场景。

\vspace{1em}
\noindent\textbf{关键词}：烟幕干扰弹；视线遮挡；遗传算法；投放策略优化；多无人机协同；运动学建模

\newpage
